\section{Conclusion}\label{sec:conclusion}

\begin{itemize}[topsep=0pt]
  \item \textbf{When.} If $A$ is TU and $b$ integral, every vertex is integral (Hoffman--Kruskal): hence integral flows, König's theorem and shift scheduling. TU is sufficient but not necessary: the VC LP was exact on 28.8\% of our non-TU graphs (E1), with half-integral vertices.
  \item \textbf{How to recognise.} A known structure can make the test trivial (bipartiteness). On cycles and several TU families exhaustive recognition became impractical around order 10--12, although small, early obstructions let it reject much larger matrices; CMR's implementation of Seymour's test decided all 340 inputs (E2).
  \item \textbf{Does the algorithm matter.} Yes. TU guarantees integral vertices, not integral optimal points: simplex and IPM with crossover returned integral answers, while IPM alone was fractional on exactly the free vertices (E3). An exact LP is not always the best tool: compiled flow and matching code was always faster, by 2--3 orders of magnitude on the largest random instances, while for shift scheduling the LP was the practical route (E4).
\end{itemize}

\paragraph{Limitations.} All LPs were solved with HiGHS, most instances were synthetic and unit-cost, our timings partly compare implementations (Python against C), and CMR's times were too close to its start-up cost to show scaling. We proved one direction of Hoffman--Kruskal and Ghouila-Houri, and saw only CMR's answers to Seymour's test, not its decomposition steps.

\paragraph{Beyond our scope.} A system $Ax\le b$ is \emph{totally dual integral} if the dual LP has an integral optimum for every integral objective with a finite optimum; with $b$ integral, it defines an integral polyhedron even when $A$ is not TU \cite{edmonds1977}. In the same spirit, Edmonds's matching polytope \cite{edmonds1965} describes matchings in non-bipartite graphs, where the incidence matrix is not TU. Current research bounds subdeterminants by a constant $\Delta$ instead of~$1$: integer programs with $\Delta=2$ (bimodular) are solvable in strongly polynomial time \cite{artmann2017}, as are $\Delta$-modular programs with at most two nonzeros per row, for fixed $\Delta$ \cite{fiorini2021}. Progress on strictly $\Delta$-modular programs continues \cite{nagele2024}, but to our knowledge it is open whether every fixed $\Delta$ admits a polynomial-time algorithm.
